\documentclass[11pt,a4paper]{article}
\usepackage[utf8]{inputenc}
\usepackage[french]{babel}
\usepackage{amsmath,amssymb,amsthm}
\usepackage{geometry}
\geometry{hmargin=2.5cm,vmargin=3cm}
\usepackage{tikz}
\usetikzlibrary{cd,positioning,shapes}
\usepackage{listings}
\usepackage{xcolor}

\lstset{
    language=Python,
    basicstyle=\ttfamily\small,
    keywordstyle=\color{blue}\bfseries,
    stringstyle=\color{red},
    commentstyle=\color{gray}\itshape,
    numbers=left,
    numberstyle=\tiny\color{gray},
    breaklines=true,
    frame=single,
    showstringspaces=false,
    literate={é}{{\'e}}1 {è}{{\`e}}1 {à}{{\`a}}1 {ê}{{\^e}}1 {î}{{\^i}}1 {ç}{{\c{c}}}1 {ï}{{\"i}}1
}

\title{\Huge JALON 53 : AXIOMES DE SÉPARATION \\ \large Cours magistral, exercices corrigés et travaux pratiques}
\author{\Large Charles EDOU NZE}
\date{\today}

\begin{document}
\maketitle
\tableofcontents
\newpage

\part{Cours Magistral}


\section*{Jalon 53 : Axiomes de séparation}

\section{Introduction aux axiomes de séparation}

Dans la théorie topologique, la notion de séparation traduit la capacité à isoler distinctement des points de l'espace à l'aide d'ouverts disjoints. Imaginez une fête où tout le monde est habillé exactement pareil. Si vous voyez deux personnes au loin, comment être sûr que ce sont bien deux individus différents et pas une illusion d'optique ? Les \textbf{Axiomes de séparation}, c'est comme une règle de politesse pour l'espace : pour que l'espace soit "propre" (Hausdorff), il faut que pour n'importe quelles deux personnes distinctes, on puisse dessiner deux cercles autour d'elles qui ne se touchent pas. Chacun a sa "bulle privée". Si l'espace ne respecte pas ça, il devient flou : deux points peuvent être si proches qu'on ne peut plus les distinguer avec des voisinages.
Sans un niveau adéquat de séparation, les notions analytiques fondamentales perdent leur robustesse : En topologie générale, certains espaces sont pathologiques. Sans axiome de séparation, une suite de nombres pourrait converger vers deux cibles différentes en même temps ! C'est catastrophique pour le calcul. On a donc défini des niveaux de "séparation" pour garantir que nos objets mathématiques se comportent de manière saine.
Géométriquement, Deux points $x$ et $y$, entourés chacun par un nuage (ouvert) tel que les deux nuages sont totalement disjoints.

\section{Définitions, Théorèmes \& Exemples Concrets}

Soit $(X, \mathcal{T})$ un espace topologique.

\subsection{L'Espace de Hausdorff ($T_2$)}

C'est l'axiome le plus important en analyse.

\begin{quote}
\textbf{Définition (Espace $T_2$ ou de Hausdorff) :}
\end{quote}
\begin{quote}
On dit que $X$ est un \textbf{espace de Hausdorff} si pour tous $x, y \in X$ tels que $x \neq y$, il existe un voisinage $U$ de $x$ et un voisinage $V$ de $y$ tels que :
\end{quote}
\begin{quote}
$U \cap V = \emptyset$
\end{quote}

\textbf{Analyse des variables et typage :}
\begin{itemize}
\item L'espace $X$ est l'ensemble de base.

\item Les points $x, y$ sont des éléments de $X$.

\item Les ensembles $U$ et $V$ sont des ouverts de la topologie $\mathcal{T}$ (donc $U, V \in \mathcal{T}$).
\end{itemize}

\textbf{Exemple Concret Numérique :}
Considérons $X = \mathbb{R}$ muni de sa topologie usuelle (définie par la métrique de la valeur absolue).
Soient $x = 1$ et $y = 2$.
Posons $r = \frac{|1 - 2|}{2} = 0.5$.
Les boules ouvertes (intervalles) correspondantes sont :
$U = B(1, 0.5) = ]0.5, 1.5[$
$V = B(2, 0.5) = ]1.5, 2.5[$
Ces deux intervalles sont ouverts dans $\mathbb{R}$, ils contiennent respectivement $1$ et $2$, et leur intersection est vide : $]0.5, 1.5[ \cap ]1.5, 2.5[ = \emptyset$. L'axiome $T_2$ est vérifié.

\textbf{Représentation Géométrique :}
\begin{center}
\begin{tikzpicture}
    \filldraw[fill=blue!10, draw=blue, thick] (0,0) circle (1.2cm);
    \filldraw[fill=red!10, draw=red, thick] (4,0) circle (1.2cm);
    \fill (0,0) circle (2pt) node[anchor=north] {$x$};
    \fill (4,0) circle (2pt) node[anchor=north] {$y$};
    \node[blue] at (-0.5, 0.5) {$U$};
    \node[red] at (4.5, 0.5) {$V$};
    \node at (2, -1.5) {L'intersection $U \cap V = \emptyset$};
\end{tikzpicture}
\end{center}


\subsection{Autres Axiomes (Pour culture)}

\begin{enumerate}
\item \textbf{$T_0$ (Kolmogorov) :} Pour deux points distincts, au moins l'un a un voisinage qui ne contient pas l'autre.

\item \textbf{$T_1$ (Fréchet) :} Chaque point a un voisinage qui ne contient pas l'autre (équivaut à dire que les singletons sont fermés).

\item \textbf{Hiérarchie :} $T_2 \implies T_1 \implies T_0$.
\end{enumerate}

\subsection{Propriétés des Espaces de Hausdorff}

\begin{quote}
\textbf{Théorème (Unicité de la limite) :}
\end{quote}
\begin{quote}
Dans un espace de Hausdorff, si une suite $(x_n)$ converge, alors sa limite est \textbf{unique}.
\end{quote}

\begin{quote}
\textbf{Théorème (Singletons) :}
\end{quote}
\begin{quote}
Dans un espace de Hausdorff, tout singleton $\{x\}$ est un ensemble \textbf{fermé}.
\end{quote}

\section{Démonstrations Rigoureuses}

\subsection{Démonstration : Unicité de la limite dans un espace $T_2$}

\begin{enumerate}
\item \textbf{Cadre :} Soit $(x_n)$ une suite dans un espace de Hausdorff $X$. Supposons par l'absurde que $x_n \to L_1$ et $x_n \to L_2$ avec $L_1 \neq L_2$.

\item \textbf{Utilisation de l'axiome $T_2$ :} Comme $L_1 \neq L_2$, il existe un ouvert $U$ contenant $L_1$ et un ouvert $V$ contenant $L_2$ tels que $U \cap V = \emptyset$.

\item \textbf{Application de la définition de limite :}
\end{enumerate}
   - Comme $x_n \to L_1$, il existe un rang $N_1$ tel que pour $n \ge N_1$, $x_n \in U$.
   - Comme $x_n \to L_2$, il existe un rang $N_2$ tel que pour $n \ge N_2$, $x_n \in V$.
\begin{enumerate}
\item \textbf{Conclusion :} Pour tout $n \ge \max(N_1, N_2)$, on a $x_n \in U \cap V$. Or $U \cap V = \emptyset$. C'est une contradiction. Donc la limite est unique.
\end{enumerate}

\subsection{Démonstration : Tout espace métrique est Hausdorff}

\begin{enumerate}
\item \textbf{Cadre :} Soit $(X, d)$ un espace métrique. Soient $x, y \in X$ avec $x \neq y$.

\item \textbf{Choix du rayon :} Posons $r = d(x, y)$. Comme $x \neq y$, $r > 0$.

\item \textbf{Construction des ouverts :} Posons $\epsilon = r/3$. Considérons les boules ouvertes $B(x, \epsilon)$ et $B(y, \epsilon)$.

\item \textbf{Vérification de la séparation :} Supposons qu'il existe $z \in B(x, \epsilon) \cap B(y, \epsilon)$.
\end{enumerate}
   Alors $d(x, z) < \epsilon$ et $d(y, z) < \epsilon$.
   Par inégalité triangulaire : $d(x, y) \le d(x, z) + d(z, y) < 2\epsilon = \frac{2}{3} r$.
\begin{enumerate}
\item \textbf{Conclusion :} On obtient $r < \frac{2}{3} r$, ce qui est impossible. Donc les boules sont disjointes. L'espace est Hausdorff.
\end{enumerate}

\section{Exercices d'Application}

\subsection{Exercice 1 : La droite avec un point doublé}
\textbf{Énoncé :} On prend deux copies de $\mathbb{R}$, notées $\mathbb{R}_a$ et $\mathbb{R}_b$. On identifie chaque $x \in \mathbb{R}_a \setminus \{0\}$ avec son correspondant dans $\mathbb{R}_b$. On obtient un espace $X$ qui ressemble à la droite réelle, mais avec "deux origines" $0_a$ et $0_b$. Cet espace est-il Hausdorff ?
\textbf{Correction Détaillée :}
Non. Tout voisinage de $0_a$ contient un intervalle du type $]-\epsilon, \epsilon[ \setminus \{0\}$. De même pour $0_b$. Ces deux voisinages se rencontreront toujours sur les points $x \neq 0$. On ne peut pas séparer les deux zéros. C'est l'exemple type d'un espace non-Hausdorff.

\subsection{Exercice 2 : Niveau Avancé (Graphe d'une fonction)}
\textbf{Énoncé :} Soit $f : X \to Y$ continue. Montrer que si $Y$ est Hausdorff, alors le graphe de $f$, $\Gamma = \{ (x, f(x)) \mid x \in X \}$, est un fermé de $X \times Y$.
\textbf{Correction Détaillée :}
On montre que le complémentaire est ouvert. Soit $(x, y) \notin \Gamma$, donc $y \neq f(x)$. Comme $Y$ est Hausdorff, il existe $V_y$ et $V_{f(x)}$ disjoints. Par continuité, $U = f^{-1}(V_{f(x)})$ est un ouvert contenant $x$. Alors $U \times V_y$ est un voisinage de $(x, y)$ qui ne rencontre pas le graphe.

\section{Application en Intelligence Artificielle}

\begin{itemize}
\item \textbf{Le Pont Théorique :} En IA, la séparation est liée à l'\textbf{Identifiabilité des paramètres}. Si l'espace des modèles n'est pas "bien séparé", on ne peut pas garantir que l'apprentissage va converger vers une solution unique et stable.

\item \textbf{Exemple Concret :}
\end{itemize}
    - \textbf{Modèles de Mélange (Gaussian Mixture Models) :} Si on ne définit pas d'ordre sur les composantes, le modèle est non-identifiable (on peut permuter les étiquettes des clusters sans changer la probabilité). L'espace des paramètres devient un espace quotient qui peut perdre certaines propriétés de séparation si on n'y prend pas garde.
    - \textbf{Embedding Disentanglement :} Dans les VAE ou les représentations auto-supervisées, on cherche à ce que chaque dimension du vecteur latent représente une caractéristique unique et "séparable" des autres. On veut éviter que deux concepts sémantiquement différents soient topologiquement inséparables dans l'espace latent.




\newpage
\part{Exercices d'Application et de Concours}


\section{Exercice 1 : Unicité de la limite dans un espace de Hausdorff \quad $\bigstar\star\star\star\star$}


\textbf{Énoncé :}
Soit $(X, \mathcal{T})$ un espace topologique. Montrer que $X$ est un espace de Hausdorff ($T_2$) si et seulement si, pour toute suite $(x_n)_{n \in \mathbb{N}}$ convergente dans $X$, sa limite est unique.

\textbf{Correction Détaillée :}
\begin{enumerate}
\item \textbf{Sens direct ($\implies$) :} Supposons $X$ Hausdorff. Soit $(x_n)$ une suite convergeant vers $L_1$ et $L_2$. Si $L_1 \neq L_2$, par séparation $T_2$, il existe des ouverts $U$ et $V$ disjoints tels que $L_1 \in U$ et $L_2 \in V$. La convergence implique qu'il existe $N_1$ tel que $\forall n \ge N_1, x_n \in U$ et $N_2$ tel que $\forall n \ge N_2, x_n \in V$. Pour $n \ge \max(N_1, N_2)$, $x_n \in U \cap V$, ce qui contredit $U \cap V = \emptyset$. Donc $L_1 = L_2$.

\item \textbf{Sens réciproque ($\impliedby$) (Remarque) :} En toute rigueur, l'implication réciproque est vraie pour les espaces à bases dénombrables de voisinages (espaces métriques par exemple). Dans un espace topologique général, l'unicité de la limite des suites ne suffit pas toujours à garantir la séparation $T_2$ (il faut utiliser des filtres ou des suites généralisées). Néanmoins, le fait que $T_2$ implique l'unicité de la limite des suites est le théorème fondamental à retenir.
\end{enumerate}





\section{Exercice 2 : La droite à deux origines \quad $\bigstar\bigstar\star\star\star$}


\textbf{Énoncé :}
On considère l'ensemble $X = (\mathbb{R} \setminus \{0\}) \cup \{0_a, 0_b\}$ muni de la topologie engendrée par la base d'ouverts suivante :
\begin{itemize}
\item Les intervalles ouverts de $\mathbb{R} \setminus \{0\}$.

\item Les ensembles de la forme $]-\epsilon, \epsilon[ \setminus \{0\} \cup \{0_a\}$ pour $\epsilon > 0$.

\item Les ensembles de la forme $]-\epsilon, \epsilon[ \setminus \{0\} \cup \{0_b\}$ pour $\epsilon > 0$.
\end{itemize}
Montrer que $X$ n'est pas un espace de Hausdorff.

\textbf{Correction Détaillée :}
Pour montrer que $X$ n'est pas Hausdorff, il suffit de trouver deux points distincts qui ne peuvent pas être séparés par des voisinages ouverts disjoints. Considérons les points $0_a$ et $0_b$.
Soit $U$ un voisinage ouvert de $0_a$ et $V$ un voisinage ouvert de $0_b$.
Par définition de la topologie sur $X$, il existe $\epsilon_1 > 0$ tel que $]-\epsilon_1, \epsilon_1[ \setminus \{0\} \cup \{0_a\} \subset U$.
De même, il existe $\epsilon_2 > 0$ tel que $]-\epsilon_2, \epsilon_2[ \setminus \{0\} \cup \{0_b\} \subset V$.
Soit $\epsilon = \min(\epsilon_1, \epsilon_2) > 0$. L'intervalle $]0, \epsilon[$ est non vide.
Soit $x \in ]0, \epsilon[$. On a $x \in ]-\epsilon_1, \epsilon_1[ \setminus \{0\} \subset U$ et $x \in ]-\epsilon_2, \epsilon_2[ \setminus \{0\} \subset V$.
Donc $x \in U \cap V$, ce qui prouve que $U \cap V \neq \emptyset$.
Les points $0_a$ et $0_b$ ne peuvent être séparés, donc $X$ n'est pas $T_2$.





\section{Exercice 3 : Sous-espaces d'un espace Hausdorff \quad $\bigstar\bigstar\star\star\star$}


\textbf{Énoncé :}
Soit $X$ un espace de Hausdorff et $Y \subset X$ muni de la topologie induite. Montrer que $Y$ est un espace de Hausdorff.

\textbf{Correction Détaillée :}
Soient $y_1, y_2 \in Y$ avec $y_1 \neq y_2$. Comme $Y \subset X$, on a aussi $y_1, y_2 \in X$.
Puisque $X$ est Hausdorff et $y_1 \neq y_2$, il existe des ouverts $U$ et $V$ de $X$ tels que $y_1 \in U$, $y_2 \in V$ et $U \cap V = \emptyset$.
Considérons les ensembles $U_Y = U \cap Y$ et $V_Y = V \cap Y$.
Par définition de la topologie induite, $U_Y$ et $V_Y$ sont des ouverts de $Y$.
De plus, $y_1 \in U_Y$ et $y_2 \in V_Y$.
Enfin, $U_Y \cap V_Y = (U \cap Y) \cap (V \cap Y) = (U \cap V) \cap Y = \emptyset \cap Y = \emptyset$.
Les points $y_1$ et $y_2$ sont séparés par des ouverts disjoints dans $Y$. Donc $Y$ est Hausdorff.





\section{Exercice 4 : Produit d'espaces Hausdorff \quad $\bigstar\bigstar\bigstar\star\star$}


\textbf{Énoncé :}
Soient $X$ et $Y$ deux espaces topologiques. Montrer que l'espace produit $X \times Y$ est Hausdorff si et seulement si $X$ et $Y$ sont Hausdorff.

\textbf{Correction Détaillée :}
\begin{enumerate}
\item \textbf{Sens direct ($\impliedby$) :} Supposons $X$ et $Y$ Hausdorff. Soient $(x_1, y_1)$ et $(x_2, y_2)$ deux points distincts de $X \times Y$.
\end{enumerate}
   - Cas 1 : $x_1 \neq x_2$. Comme $X$ est Hausdorff, il existe $U_1, U_2$ ouverts de $X$ disjoints contenant respectivement $x_1$ et $x_2$. Alors $U_1 \times Y$ et $U_2 \times Y$ sont des ouverts de $X \times Y$, disjoints, contenant respectivement $(x_1, y_1)$ et $(x_2, y_2)$.
   - Cas 2 : $y_1 \neq y_2$. Symétriquement, il existe $V_1, V_2$ ouverts de $Y$ disjoints. $X \times V_1$ et $X \times V_2$ séparent les points.
   Dans tous les cas, $X \times Y$ est Hausdorff.
\begin{enumerate}
\item \textbf{Sens réciproque ($\implies$) :} Supposons $X \times Y$ Hausdorff. Soient $x_1, x_2 \in X$ avec $x_1 \neq x_2$. Fixons $y \in Y$. Les points $(x_1, y)$ et $(x_2, y)$ sont distincts dans $X \times Y$. Ils admettent des voisinages de base disjoints $U_1 \times V_1$ et $U_2 \times V_2$. L'intersection est vide si et seulement si $U_1 \cap U_2 = \emptyset$ ou $V_1 \cap V_2 = \emptyset$. Comme $y \in V_1 \cap V_2$, on a nécessairement $U_1 \cap U_2 = \emptyset$. Les ouverts $U_1$ et $U_2$ de $X$ séparent $x_1$ et $x_2$. Donc $X$ est Hausdorff (même argument pour $Y$).
\end{enumerate}





\section{Exercice 5 : Fermeture de la diagonale \quad $\bigstar\bigstar\bigstar\star\star$}


\textbf{Énoncé :}
Soit $X$ un espace topologique. Montrer que $X$ est Hausdorff si et seulement si la diagonale $\Delta = \{(x, x) \mid x \in X\}$ est fermée dans l'espace produit $X \times X$.

\textbf{Correction Détaillée :}
\begin{enumerate}
\item \textbf{Sens direct ($\implies$) :} Supposons $X$ Hausdorff. Pour montrer que $\Delta$ est fermée, on montre que son complémentaire $(X \times X) \setminus \Delta$ est ouvert. Soit $(x, y) \notin \Delta$, i.e., $x \neq y$. Par hypothèse, il existe $U$ ouvert contenant $x$ et $V$ ouvert contenant $y$ tels que $U \cap V = \emptyset$. L'ensemble $U \times V$ est un ouvert de $X \times X$ contenant $(x, y)$. De plus, $(U \times V) \cap \Delta = \emptyset$ car si $(z, z) \in U \times V$, alors $z \in U \cap V$, ce qui est impossible. Ainsi $(x, y) \in U \times V \subset (X \times X) \setminus \Delta$. Le complémentaire est ouvert, $\Delta$ est fermée.

\item \textbf{Sens réciproque ($\impliedby$) :} Supposons $\Delta$ fermée. Soient $x \neq y$ dans $X$. Alors $(x, y) \in (X \times X) \setminus \Delta$, qui est ouvert. Par définition de la topologie produit, il existe un ouvert de base $U \times V$ tel que $(x, y) \in U \times V \subset (X \times X) \setminus \Delta$. On a $x \in U$, $y \in V$. Si $U \cap V$ contenait un élément $z$, alors $(z, z) \in U \times V$, ce qui contredirait l'inclusion dans le complémentaire de $\Delta$. Donc $U \cap V = \emptyset$. $X$ est Hausdorff.
\end{enumerate}





\section{Exercice 6 : Espace de Zariski \quad $\bigstar\bigstar\bigstar\bigstar\star$}


\textbf{Énoncé :}
Soit $X$ un ensemble infini muni de la topologie cofinie (un sous-ensemble est ouvert si et seulement s'il est vide ou si son complémentaire est fini). Cet espace est-il Hausdorff ? Est-il $T_1$ (axiome de Fréchet) ?

\textbf{Correction Détaillée :}
\begin{enumerate}
\item \textbf{Axiome $T_1$ :} Soit $x \in X$. Le complémentaire du singleton $\{x\}$ est $X \setminus \{x\}$. Comme $\{x\}$ est fini (cardinal 1), $X \setminus \{x\}$ est ouvert, donc $\{x\}$ est fermé. Un espace dont tous les singletons sont fermés est un espace $T_1$. Donc oui, $X$ est $T_1$.

\item \textbf{Axiome $T_2$ (Hausdorff) :} Soient $x, y \in X$ avec $x \neq y$. Supposons qu'il existe des ouverts $U$ et $V$ disjoints tels que $x \in U$ et $y \in V$. Comme $U$ et $V$ sont non vides, leurs complémentaires $U^c$ et $V^c$ sont finis. Or, $U \cap V = \emptyset \implies (U \cap V)^c = X \implies U^c \cup V^c = X$. L'union de deux ensembles finis étant finie, $X$ serait fini, ce qui contredit l'hypothèse. Il est donc impossible de trouver de tels ouverts. $X$ n'est pas Hausdorff.
\end{enumerate}





\section{Exercice 7 : Graphe d'une fonction continue \quad $\bigstar\bigstar\bigstar\bigstar\star$}


\textbf{Énoncé :}
Soit $f : X \to Y$ une application continue entre deux espaces topologiques. On suppose que $Y$ est Hausdorff. Montrer que le graphe de $f$, défini par $\Gamma = \{(x, f(x)) \mid x \in X\}$, est fermé dans $X \times Y$.

\textbf{Correction Détaillée :}
On va montrer que le complémentaire de $\Gamma$ est ouvert.
Soit $(x, y) \in (X \times Y) \setminus \Gamma$. Cela signifie que $y \neq f(x)$.
Puisque $Y$ est Hausdorff, il existe des ouverts disjoints $V_1, V_2 \subset Y$ tels que $f(x) \in V_1$ et $y \in V_2$.
Comme $f$ est continue, $U = f^{-1}(V_1)$ est un ouvert de $X$ contenant $x$.
Considérons l'ouvert $U \times V_2$ de $X \times Y$. Ce voisinage contient $(x, y)$.
De plus, si $(x', y') \in U \times V_2$, alors $x' \in U \implies f(x') \in V_1$. Et $y' \in V_2$.
Puisque $V_1 \cap V_2 = \emptyset$, on a $y' \neq f(x')$, ce qui signifie que $(x', y') \notin \Gamma$.
Ainsi, $U \times V_2 \subset (X \times Y) \setminus \Gamma$. Le complémentaire est ouvert, donc $\Gamma$ est fermé.





\section{Exercice 8 : Espaces normaux ($T_4$) et lemmes de séparation \quad $\bigstar\bigstar\bigstar\bigstar\star$}


\textbf{Énoncé :}
Rappeler la définition d'un espace normal ($T_4$). Montrer que tout espace métrique est normal.

\textbf{Correction Détaillée :}
\begin{enumerate}
\item \textbf{Définition :} Un espace topologique $X$ est normal ($T_4$) s'il est $T_1$ (les singletons sont fermés) et si pour tout couple $(A, B)$ de fermés disjoints, il existe des ouverts $U, V$ disjoints tels que $A \subset U$ et $B \subset V$.

\item \textbf{Cas métrique :} Soit $(X, d)$ un espace métrique. Les singletons sont fermés car l'espace est Hausdorff (et même métrique). Soient $A, B$ deux fermés disjoints. Pour tout $a \in A$, $a \notin B$. $B$ étant fermé, $d(a, B) = \inf_{b \in B} d(a, b) > 0$. Posons $r_a = \frac{1}{2} d(a, B)$. L'ouvert $U = \bigcup_{a \in A} B(a, r_a)$ contient $A$. Symétriquement, pour $b \in B$, posons $r_b = \frac{1}{2} d(b, A)$ et $V = \bigcup_{b \in B} B(b, r_b)$. L'ouvert $V$ contient $B$.
\end{enumerate}
Montrons $U \cap V = \emptyset$. Si $x \in U \cap V$, il existe $a \in A$ et $b \in B$ tels que $d(x, a) < r_a$ et $d(x, b) < r_b$. Par inégalité triangulaire, $d(a, b) \le d(a, x) + d(x, b) < r_a + r_b$.
Or, $d(a, B) \le d(a, b)$ et $d(b, A) \le d(a, b)$. Donc $2r_a \le d(a, b)$ et $2r_b \le d(a, b)$. Ainsi $r_a + r_b \le \max(2r_a, 2r_b) \le d(a, b)$. Contradiction avec $d(a, b) < r_a + r_b$.
Donc $U \cap V = \emptyset$. L'espace métrique est normal.





\section{Exercice 9 : Espace quotient et séparation \quad $\bigstar\bigstar\bigstar\bigstar\bigstar$}


\textbf{Énoncé :}
Soit $X = \mathbb{R}$ et la relation d'équivalence $x \sim y \iff x - y \in \mathbb{Q}$. On munit l'espace quotient $X / \sim$ de la topologie quotient. Montrer que $X / \sim$ est un espace topologique grossier (indiscret), et en déduire qu'il n'est pas Hausdorff.

\textbf{Correction Détaillée :}
Soit $\pi : \mathbb{R} \to \mathbb{R}/\mathbb{Q}$ la projection canonique. Par définition, $U \subset \mathbb{R}/\mathbb{Q}$ est ouvert si et seulement si $\pi^{-1}(U)$ est ouvert dans $\mathbb{R}$.
Supposons qu'il existe un ouvert non vide $U$ dans $X / \sim$ tel que $U \neq X / \sim$.
Alors $\pi^{-1}(U)$ est un ouvert non vide de $\mathbb{R}$ tel que $\pi^{-1}(U) \neq \mathbb{R}$.
Puisque $U$ est non vide, il existe $\bar{x} \in U$, donc $\pi^{-1}(\{\bar{x}\}) \subset \pi^{-1}(U)$.
Or, $\pi^{-1}(\{\bar{x}\}) = x + \mathbb{Q}$, qui est dense dans $\mathbb{R}$.
Un ouvert $\pi^{-1}(U)$ contenant une partie dense est nécessairement dense. Et un ouvert n'est égal à $\mathbb{R}$ que s'il est vide ou plein ? Non. L'ouvert $\pi^{-1}(U)$ a la propriété supplémentaire d'être saturé pour la relation, c'est-à-dire que si $y \in \pi^{-1}(U)$, alors $y + \mathbb{Q} \subset \pi^{-1}(U)$.
Comme $\pi^{-1}(U)$ contient au moins un point $x$, il contient $x + \mathbb{Q}$. Comme $\pi^{-1}(U)$ est ouvert, il contient un intervalle $]a, b[$. La réunion des translatés rationnels de cet intervalle est $\mathbb{R}$ tout entier : $\bigcup_{q \in \mathbb{Q}} (]a, b[ + q) = \mathbb{R}$.
Donc $\pi^{-1}(U) = \mathbb{R}$, d'où $U = X / \sim$.
Les seuls ouverts sont $\emptyset$ et $X / \sim$. L'espace est grossier.
Puisqu'il a plus d'un point et que les seuls ouverts sont triviaux, il est impossible de séparer deux points. Il n'est pas Hausdorff.





\section{Exercice 10 : Compacité et séparation \quad $\bigstar\bigstar\bigstar\bigstar\bigstar$}


\textbf{Énoncé :}
Montrer que toute application continue et bijective d'un espace topologique compact $X$ vers un espace topologique Hausdorff $Y$ est un homéomorphisme.

\textbf{Correction Détaillée :}
Soit $f : X \to Y$ continue et bijective. Il faut montrer que $f^{-1} : Y \to X$ est continue, ce qui équivaut à montrer que pour tout fermé $F$ de $X$, $(f^{-1})^{-1}(F) = f(F)$ est fermé dans $Y$.
Soit $F$ un fermé de $X$. Comme $X$ est compact et que tout fermé d'un compact est compact, $F$ est compact.
L'image continue d'un compact est compacte. Donc $f(F)$ est compact dans $Y$.
Puisque $Y$ est un espace Hausdorff, et que tout sous-espace compact d'un espace Hausdorff est fermé, on conclut que $f(F)$ est fermé.
Ainsi, $f$ est une application fermée, et bijective. Donc $f^{-1}$ est continue. $f$ est un homéomorphisme.
C'est un théorème central liant compacité et séparation.




\newpage
\part{Travaux Pratiques et Simulations Algorithmiques}


\subsection{TP 1 : Implémentation d'une topologie cofinie}


Dans ce TP, nous modélisons une topologie cofinie sur un ensemble fini, et nous vérifions si elle respecte l'axiome de séparation de Kolmogorov (T0) et de Fréchet (T1).
L'ensemble $X$ est simulé par une liste d'entiers.

\begin{lstlisting}
def is_cofinite_topology(X, open_sets):
    # Verifie si open_sets forme une topologie cofinie sur X
    # Un ensemble est ouvert s'il est vide ou si son complementaire est fini.
    # Puisque X est fini, tous les sous-ensembles devraient etre ouverts
    # pour que ce soit LA topologie cofinie, mais dans l'esprit general,
    # c'est la topologie discrete pour X fini.
    # Simulons plutot une topologie specifique pour tester T0 et T1.
    pass

def is_T0(X, open_sets):
    for x in X:
        for y in X:
            if x != y:
                # Chercher un ouvert contenant x et pas y, OU contenant y et pas x
                separated = False
                for u in open_sets:
                    if (x in u and y not in u) or (y in u and x not in u):
                        separated = True
                        break
                if not separated:
                    return False
    return True

def is_T1(X, open_sets):
    for x in X:
        for y in X:
            if x != y:
                # Chercher un ouvert contenant x et pas y, ET un ouvert contenant y et pas x
                x_sep = False
                y_sep = False
                for u in open_sets:
                    if x in u and y not in u:
                        x_sep = True
                    if y in u and x not in u:
                        y_sep = True
                if not (x_sep and y_sep):
                    return False
    return True

X = [1, 2, 3]
# Topologie de Sierpinski : X, vide, {1}
open_sets_sierpinski = [[], [1], [1, 2, 3]]
assert is_T0(X, open_sets_sierpinski) == False # 2 et 3 ne sont pas separes du tout
# En fait, 1 et 2 sont T0, mais 2 et 3 partagent exactement les memes ouverts.

X_sierp = [1, 2]
open_sets_sierp_2 = [[], [1], [1, 2]]
assert is_T0(X_sierp, open_sets_sierp_2) == True
assert is_T1(X_sierp, open_sets_sierp_2) == False # 2 ne peut pas etre isole de 1

open_sets_discrete = [[], [1], [2], [1, 2]]
assert is_T1(X_sierp, open_sets_discrete) == True

print("Tests T0/T1 passes.")
\end{lstlisting}





\subsection{TP 2 : Vérification de l'axiome de Hausdorff (T2)}


Nous continuons avec notre modélisation d'espaces topologiques finis. Cette fois, nous implémentons un test pour vérifier si la topologie donnée est Hausdorff (T2).

\begin{lstlisting}
def is_T2_hausdorff(X, open_sets):
    for i in range(len(X)):
        for j in range(i + 1, len(X)):
            x = X[i]
            y = X[j]
            separated = False
            for u in open_sets:
                for v in open_sets:
                    if x in u and y not in u and y in v and x not in v:
                        # Verifier l'intersection vide
                        intersection_empty = True
                        for elem in u:
                            if elem in v:
                                intersection_empty = False
                                break
                        if intersection_empty:
                            separated = True
                            break
                if separated:
                    break
            if not separated:
                return False
    return True

X = [1, 2, 3]
open_sets_discrete = [[], [1], [2], [3], [1, 2], [1, 3], [2, 3], [1, 2, 3]]
assert is_T2_hausdorff(X, open_sets_discrete) == True

open_sets_not_t2 = [[], [1, 2], [3], [1, 2, 3]]
assert is_T2_hausdorff(X, open_sets_not_t2) == False # 1 et 2 ne peuvent pas etre separes

print("Tests Hausdorff (T2) passes.")
\end{lstlisting}





\subsection{TP 3 : La droite avec point double en Python}


Nous allons simuler l'espace de la droite avec deux origines en utilisant des voisinages.
Nous utiliserons des intervalles discrets pour représenter les ouverts de base et démontrer programmatiquement que les deux origines ne peuvent être séparées.

\begin{lstlisting}
def generate_neighborhoods_0a(epsilon):
    return set(range(-epsilon, epsilon + 1)) - {0} | {'0a'}

def generate_neighborhoods_0b(epsilon):
    return set(range(-epsilon, epsilon + 1)) - {0} | {'0b'}

# On tente de trouver epsilon1 et epsilon2 tels que les voisinages sont disjoints
def are_0a_0b_separable(max_epsilon):
    for eps1 in range(1, max_epsilon + 1):
        for eps2 in range(1, max_epsilon + 1):
            V_0a = generate_neighborhoods_0a(eps1)
            V_0b = generate_neighborhoods_0b(eps2)

            # Intersection
            intersection = V_0a.intersection(V_0b)
            if len(intersection) == 0:
                return True # On a reussi a les separer

    return False # Impossible a separer pour les epsilons testes

assert are_0a_0b_separable(10) == False

print("Simulation de la droite a deux origines reussie, pas de separation.")
\end{lstlisting}





\subsection{TP 4 : Simulation du lemme d'Urysohn discret}


Un espace normal (T4) vérifie le lemme d'Urysohn : on peut séparer des fermés disjoints par une fonction continue.
Nous allons implémenter un mini-solveur pour un espace fini qui construit cette fonction séparatrice (valant 0 sur A et 1 sur B).

\begin{lstlisting}
def check_normal(X, closed_sets):
    # Pour tout A, B fermes disjoints, il faut U, V ouverts disjoints avec A subset U, B subset V
    open_sets = []
    for c in closed_sets:
        open_sets.append([x for x in X if x not in c])

    for A in closed_sets:
        for B in closed_sets:
            # Check si disjoints
            disjoint = True
            for a in A:
                if a in B:
                    disjoint = False

            if disjoint and len(A) > 0 and len(B) > 0:
                separated = False
                for U in open_sets:
                    for V in open_sets:
                        # A dans U ?
                        A_in_U = all(a in U for a in A)
                        # B dans V ?
                        B_in_V = all(b in V for b in B)

                        if A_in_U and B_in_V:
                            U_V_disjoint = True
                            for u_elem in U:
                                if u_elem in V:
                                    U_V_disjoint = False
                            if U_V_disjoint:
                                separated = True

                if not separated:
                    return False
    return True

X = [1, 2, 3, 4]
# Topologie grossiere (T4 trivialement car pas de fermes disjoints non vides)
closed_grossiere = [[], [1, 2, 3, 4]]
assert check_normal(X, closed_grossiere) == True

# Une topologie discrete est normale
# Generation de toutes les parties
from itertools import chain, combinations
def powerset(iterable):
    s = list(iterable)
    return list(chain.from_iterable(combinations(s, r) for r in range(len(s)+1)))

closed_discrete = [list(c) for c in powerset(X)]
assert check_normal(X, closed_discrete) == True

print("Tests de normalite T4 passes.")
\end{lstlisting}





\subsection{TP 5 : Propriété de séparation et graphe fermé}


Dans ce TP, nous simulons une fonction $f: X \to Y$. Nous vérifions que si l'espace $Y$ a une "diagonale fermée" (i.e. T2), alors le graphe de la fonction continue est fermé dans l'espace produit.
Nous utilisons des graphes et des listes d'adjacence pour modéliser la proximité.

\begin{lstlisting}
def is_diagonal_closed(Y):
    # Simulation: dans un espace discret fini, la diagonale est toujours fermee.
    diagonal = [(y, y) for y in Y]
    # On simule la fermeture en verifiant que pour tout element hors diagonale,
    # on peut trouver un voisinage produit qui ne touche pas la diagonale.
    for y1 in Y:
        for y2 in Y:
            if y1 != y2:
                # Voisinages discrets : {y1} et {y2}
                # Produit : {(y1, y2)} inter diagonal est vide
                pass
    return True

def graph_is_closed(X, Y, f_dict):
    # X et Y discrets
    graph = [(x, f_dict[x]) for x in X]

    # On verifie que le complementaire est ouvert (donc union de ouverts de base)
    # Dans un espace discret, tout est ouvert, donc le graphe est ferme.
    # Ceci est une verification de consistance triviale.
    for x in X:
        for y in Y:
            if (x, y) not in graph:
                # (x, y) est dans le complementaire
                assert y != f_dict[x]
    return True

X = [1, 2, 3]
Y = ['a', 'b', 'c']
f = {1: 'a', 2: 'b', 3: 'a'}

assert is_diagonal_closed(Y) == True
assert graph_is_closed(X, Y, f) == True

print("Tests de graphe ferme passes.")
\end{lstlisting}



\end{document}